%% style/macros.tex
%% Math notation macros, structural macros (\whereWeAre, \chapterRecap), and Colab links.

%%%%%%%%%%%%%%%%%%%%%%%%%%%%%%%%%%%%%%%%%%%%%%%%%%%%%%%%%%%%%%%%%%%%%%%%%%%%%%%
% Number sets and common math objects
%%%%%%%%%%%%%%%%%%%%%%%%%%%%%%%%%%%%%%%%%%%%%%%%%%%%%%%%%%%%%%%%%%%%%%%%%%%%%%%

\providecommand{\R}{\mathbb{R}}
\providecommand{\N}{\mathbb{N}}
\providecommand{\Z}{\mathbb{Z}}
\providecommand{\Q}{\mathbb{Q}}
\providecommand{\C}{\mathbb{C}}

\providecommand{\E}{\mathbb{E}}
\providecommand{\Prob}{\mathbb{P}}
\providecommand{\Var}{\mathrm{Var}}
\providecommand{\Cov}{\mathrm{Cov}}
\providecommand{\Cor}{\mathrm{Cor}}

% Common distributions
\providecommand{\Normal}{\mathcal{N}}
\providecommand{\Uniform}{\mathcal{U}}
\providecommand{\Wishart}{\mathcal{W}}

% Convergence arrows
\providecommand{\convas}{\xrightarrow{\mathrm{a.s.}}}
\providecommand{\convp}{\xrightarrow{P}}
\providecommand{\convd}{\xrightarrow{d}}

% Linear algebra
\providecommand{\transpose}{^{\mathsf{T}}}
\providecommand{\inv}{^{-1}}
\providecommand{\Tr}{\mathrm{Tr}}
\providecommand{\rank}{\mathrm{rank}}
\providecommand{\diag}{\mathrm{diag}}

% Sample covariance
\providecommand{\Shat}{\widehat{\Sigma}_n}
\providecommand{\Spop}{\Sigma}

% Vectors and matrices: light convention, bold for vectors, capital for matrices
\providecommand{\vect}[1]{\boldsymbol{#1}}
\providecommand{\mat}[1]{\boldsymbol{#1}}

% Useful operators
\DeclareMathOperator*{\argmin}{arg\,min}
\DeclareMathOperator*{\argmax}{arg\,max}
\DeclareMathOperator{\sgn}{sgn}

%%%%%%%%%%%%%%%%%%%%%%%%%%%%%%%%%%%%%%%%%%%%%%%%%%%%%%%%%%%%%%%%%%%%%%%%%%%%%%%
% Vocabulary indexing: \vocab{spectral theorem} bolds and indexes the term
%%%%%%%%%%%%%%%%%%%%%%%%%%%%%%%%%%%%%%%%%%%%%%%%%%%%%%%%%%%%%%%%%%%%%%%%%%%%%%%

\providecommand{\vocab}[1]{\index{#1}\emph{#1}}
\providecommand{\vocabalt}[2]{\index{#2}\emph{#1}}

%%%%%%%%%%%%%%%%%%%%%%%%%%%%%%%%%%%%%%%%%%%%%%%%%%%%%%%%%%%%%%%%%%%%%%%%%%%%%%%
% Colab linking, every \colablink{NAME}{URL} produces a clickable badge
%%%%%%%%%%%%%%%%%%%%%%%%%%%%%%%%%%%%%%%%%%%%%%%%%%%%%%%%%%%%%%%%%%%%%%%%%%%%%%%

\newcommand{\colablink}[2]{%
  \href{#2}{%
    \begin{tikzpicture}[baseline=(badge.base)]
      \node[badge/.style={rectangle, rounded corners=2pt, draw=mathemagicsColab,
            fill=mathemagicsColab!20, inner sep=2pt, font=\scriptsize\sffamily}]
        (badge) {Colab: #1};
    \end{tikzpicture}%
  }%
}

%%%%%%%%%%%%%%%%%%%%%%%%%%%%%%%%%%%%%%%%%%%%%%%%%%%%%%%%%%%%%%%%%%%%%%%%%%%%%%%
% \whereWeAre{...}, half-page mini-map at the start of each chapter
% Place at the top of the chapter, after \chapter{...}
%%%%%%%%%%%%%%%%%%%%%%%%%%%%%%%%%%%%%%%%%%%%%%%%%%%%%%%%%%%%%%%%%%%%%%%%%%%%%%%

\newtcolorbox{whereWeAreBox}{%
  enhanced,
  colback=mathemagicsWhereWeAreBg,
  colframe=mathemagicsWhereWeAre,
  boxrule=0.4pt,
  arc=2pt,
  fonttitle=\bfseries\sffamily\small,
  title={Where we are},
  attach boxed title to top left={yshift=-2mm, xshift=4mm},
  boxed title style={colback=mathemagicsWhereWeAre, colframe=mathemagicsWhereWeAre,
                     sharp corners, boxrule=0pt},
  coltitle=white,
  top=4mm,
  before skip=0.6em,
  after skip=1.2em,
}

\newcommand{\whereWeAre}[1]{%
  \begin{whereWeAreBox}%
    #1%
  \end{whereWeAreBox}%
}

%%%%%%%%%%%%%%%%%%%%%%%%%%%%%%%%%%%%%%%%%%%%%%%%%%%%%%%%%%%%%%%%%%%%%%%%%%%%%%%
% \chapterRecap{...}, last page of each chapter, 3-5 punchline bullets
%%%%%%%%%%%%%%%%%%%%%%%%%%%%%%%%%%%%%%%%%%%%%%%%%%%%%%%%%%%%%%%%%%%%%%%%%%%%%%%

\newtcolorbox{chapterRecapBox}{%
  enhanced,
  colback=mathemagicsRecapBg,
  colframe=mathemagicsRecap,
  boxrule=0.4pt,
  arc=2pt,
  fonttitle=\bfseries\sffamily\small,
  title={Chapter recap},
  attach boxed title to top left={yshift=-2mm, xshift=4mm},
  boxed title style={colback=mathemagicsRecap, colframe=mathemagicsRecap,
                     sharp corners, boxrule=0pt},
  coltitle=white,
  top=4mm,
  before skip=1.2em,
  after skip=0.6em,
}

\newcommand{\chapterRecap}[1]{%
  \begin{chapterRecapBox}%
    #1%
  \end{chapterRecapBox}%
}

%%%%%%%%%%%%%%%%%%%%%%%%%%%%%%%%%%%%%%%%%%%%%%%%%%%%%%%%%%%%%%%%%%%%%%%%%%%%%%%
% Difficulty pictogram, \coffee{1}, \coffee{2}, \coffee{3} for exercise rating
%%%%%%%%%%%%%%%%%%%%%%%%%%%%%%%%%%%%%%%%%%%%%%%%%%%%%%%%%%%%%%%%%%%%%%%%%%%%%%%

\newcommand{\coffee}[1]{%
  {\color{mathemagicsExercise}%
  \ifcase#1\or\textbf{$\bullet$}\or\textbf{$\bullet\bullet$}\or\textbf{$\bullet\bullet\bullet$}\fi}%
}
% Note: this is a placeholder. Step 4 of the plan replaces with actual coffee-cup glyph
% (a small TikZ icon) once the visual identity is locked.

%%%%%%%%%%%%%%%%%%%%%%%%%%%%%%%%%%%%%%%%%%%%%%%%%%%%%%%%%%%%%%%%%%%%%%%%%%%%%%%
% Cartoon-aside pictogram, used to flag non-load-bearing margin asides
% Currently a simple star; will be replaced with custom TikZ glyph in Step 4.
%%%%%%%%%%%%%%%%%%%%%%%%%%%%%%%%%%%%%%%%%%%%%%%%%%%%%%%%%%%%%%%%%%%%%%%%%%%%%%%

\newcommand{\cartoonglyph}{{\color{mathemagicsCartoon}\ensuremath{\star}}}

%%%%%%%%%%%%%%%%%%%%%%%%%%%%%%%%%%%%%%%%%%%%%%%%%%%%%%%%%%%%%%%%%%%%%%%%%%%%%%%
% Physical-object icons (in-house TikZ).
% These are inline-sized so they can sit in prose, in margins, or in tables.
% Usage: \coin{H}, \coin{T}, \die{1}, \die{2}, ..., \die{6}
% Bigger versions for figures: \coinbig{H}, \diebig{3}.
%%%%%%%%%%%%%%%%%%%%%%%%%%%%%%%%%%%%%%%%%%%%%%%%%%%%%%%%%%%%%%%%%%%%%%%%%%%%%%%

% Small inline coin (~1.1ex tall, sits on the baseline)
\DeclareRobustCommand{\coin}[1]{%
  \tikz[baseline=-0.4ex]{%
    \filldraw[draw=mathemagicsThm, line width=0.35pt,
              fill=mathemagicsIntuBg!70]
      (0,0) circle (0.55ex);
    \node[font=\tiny\sffamily\bfseries, color=mathemagicsThm,
          inner sep=0pt, outer sep=0pt] at (0,0.05ex) {#1};
  }%
}

% Larger coin for use in figures (~3ex tall)
\DeclareRobustCommand{\coinbig}[1]{%
  \tikz[baseline=-0.5ex]{%
    \filldraw[draw=mathemagicsThm, line width=0.5pt,
              fill=mathemagicsIntuBg!70]
      (0,0) circle (1.5ex);
    \node[font=\small\sffamily\bfseries, color=mathemagicsThm,
          inner sep=0pt, outer sep=0pt] at (0,0) {#1};
  }%
}

% Helper: render the pip pattern for face N at the given pip-size.
% Implemented with \ifnum chains (more robust than \ifcase under expansion).
\newcommand{\mathemagics@diepips}[2]{%
  % #1 = face number (1..6)
  % #2 = pip radius (in current units)
  \ifnum#1=1 \fill[mathemagicsThm] (0,0) circle (#2);\fi
  \ifnum#1=2 \fill[mathemagicsThm] (-0.25,0.25) circle (#2);
             \fill[mathemagicsThm] (0.25,-0.25) circle (#2);\fi
  \ifnum#1=3 \fill[mathemagicsThm] (-0.25,0.25) circle (#2);
             \fill[mathemagicsThm] (0,0) circle (#2);
             \fill[mathemagicsThm] (0.25,-0.25) circle (#2);\fi
  \ifnum#1=4 \fill[mathemagicsThm] (-0.25,-0.25) circle (#2);
             \fill[mathemagicsThm] (-0.25,0.25) circle (#2);
             \fill[mathemagicsThm] (0.25,-0.25) circle (#2);
             \fill[mathemagicsThm] (0.25,0.25) circle (#2);\fi
  \ifnum#1=5 \fill[mathemagicsThm] (-0.25,-0.25) circle (#2);
             \fill[mathemagicsThm] (-0.25,0.25) circle (#2);
             \fill[mathemagicsThm] (0,0) circle (#2);
             \fill[mathemagicsThm] (0.25,-0.25) circle (#2);
             \fill[mathemagicsThm] (0.25,0.25) circle (#2);\fi
  \ifnum#1=6 \fill[mathemagicsThm] (-0.25,-0.30) circle (#2);
             \fill[mathemagicsThm] (-0.25,0) circle (#2);
             \fill[mathemagicsThm] (-0.25,0.30) circle (#2);
             \fill[mathemagicsThm] (0.25,-0.30) circle (#2);
             \fill[mathemagicsThm] (0.25,0) circle (#2);
             \fill[mathemagicsThm] (0.25,0.30) circle (#2);\fi
}

% Small inline die (~1.4ex tall) with proper pip pattern
\DeclareRobustCommand{\die}[1]{%
  \begin{tikzpicture}[baseline=-0.5ex, x=1.6ex, y=1.6ex]
    \filldraw[rounded corners=0.05, draw=mathemagicsThm, line width=0.35pt,
              fill=mathemagicsExBg]
      (-0.5,-0.5) rectangle (0.5,0.5);
    \mathemagics@diepips{#1}{0.07}
  \end{tikzpicture}%
}

% Larger die for figures (~3ex tall)
\DeclareRobustCommand{\diebig}[1]{%
  \begin{tikzpicture}[baseline=-0.5ex, x=3ex, y=3ex]
    \filldraw[rounded corners=0.08, draw=mathemagicsThm, line width=0.5pt,
              fill=mathemagicsExBg]
      (-0.5,-0.5) rectangle (0.5,0.5);
    \mathemagics@diepips{#1}{0.08}
  \end{tikzpicture}%
}
